% engines: lualatex, pdflatex
% expect-warnings: 0
% expect-errors: 0
\documentclass{article}
\usepackage[margin=1cm]{geometry}
\usepackage{../../latex/pgf-go}
\newcommand\case[1]{\par\medskip\noindent\textbf{#1}\par\nopagebreak}
\pagestyle{empty}
\begin{document}

\case{partial = auto: a corner, with one more line; the sides near the edge reach it}
\begin{goban}[19, partial = auto, label = {at = all}]
	\stones[b]{Q16, R14}\stones[w]{R16, P17}
\end{goban}
\begin{goban}[19, partial = auto, label = {at = all}]
	\moves{Q16, R14, R16, P17}
\end{goban}

\case{partial margin = 0 and 3; in the middle of the board}
\begin{goban}[19, partial = auto, partial margin = 0, label = {at = all}]
	\stones[b]{J10, K11}\stones[w]{K9}
\end{goban}
\begin{goban}[19, partial = auto, partial margin = 3, label = {at = all}]
	\stones[b]{J10, K11}\stones[w]{K9}
\end{goban}

\case{Marks, arrows and territory are part of what is drawn}
\begin{goban}[19, partial = auto]
	\stones[b]{D4}\marks[triangle]{F7}
\end{goban}
\begin{goban}[19, partial = auto]
	\stones[b]{D4}\arrow{D4}{H6}
\end{goban}
\begin{goban}[19, partial = auto]
	\stones[b]{C3, D3, D2}\territory[b, draw = region]{C2, C1, D1}\territory[w, draw = fill]{H8}
\end{goban}

\case{An empty board is whole}
\begin{goban}[9, partial = auto]\end{goban}

\case{partial = auto and then false, or a window: the last one wins}
\begin{goban}[9, partial = auto, partial = false]\stones[b]{C3}\end{goban}
\begin{goban}[9, partial = auto, partial = A1:E5]\stones[b]{C3, G7}\end{goban}

\end{document}
